% conclusion.tex -- BNR-2026-012

\section{Conclusion}
\label{sec:conclusion}

This paper presented a production Data Availability Committee
architecture for the Basis Network enterprise zkEVM L2. The hybrid
AES-256-GCM + Reed-Solomon $(5,7)$ + Shamir $(5,7)$-SS protocol
achieves computational data privacy with information-theoretic key
secrecy, matching the security requirements of enterprise deployments
while delivering substantial performance improvements over the pure
Shamir baseline (RU-V6): $36\times$ faster attestation, $2{,}613\times$
faster recovery, and $2.77\times$ lower storage overhead.

The key architectural insight is the separation of concerns between
privacy (AES encryption) and redundancy (RS erasure coding), with Shamir
secret sharing applied only to the 32-byte AES key. This decomposition
eliminates the $O(k^2)$ per-element cost of full Shamir sharing while
retaining information-theoretic secrecy for the most critical secret
(the decryption key).

All four hypothesis components were confirmed with significant margins:
availability exceeds 99.997\% at $p = 0.99$ (target: 99.9\%);
attestation completes in 8.9\,ms at 1\,MB (target: $<$1\,s); recovery
achieves 100\% success with byte-perfect reconstruction; storage
overhead matches the theoretical minimum of $1.40\times$.

Six invariants (INV-DA-P1 through INV-DA-P6) formalize the protocol's
guarantees for downstream TLA+ specification and Coq verification.
The invariants cover verifiable encoding, data recoverability,
enterprise privacy, attestation liveness, fallback safety, and
availability guarantees.

\subsection{Recommendations}

Based on these results, we recommend the following for the Basis Network
production deployment:

\begin{enumerate}
  \item Adopt the hybrid AES+RS+Shamir architecture as the canonical
    DAC protocol for all enterprise L2 chains.
  \item Deploy 7-node committees with $k = 5$ threshold, providing
    2-node fault tolerance and 3-node honest minority for attestation
    blocking.
  \item Implement BLS signature aggregation on BN254 for constant-size
    on-chain attestation (48 bytes).
  \item Add KZG polynomial commitments for verifiable encoding,
    leveraging the BN128 pairing precompiles available on Subnet-EVM.
  \item Implement AnyTrust-style fallback to on-chain DA when fewer than
    $k$ nodes are available.
\end{enumerate}

\subsection{Future Work}

Immediate next steps in the R\&D pipeline:

\begin{itemize}
  \item \textbf{Stage 2}: Network latency simulation (50--100\,ms RTT)
    and KZG commitment generation benchmarks.
  \item \textbf{Stage 3}: Adversarial scenarios---malicious node sending
    corrupt chunks, Byzantine behavior, Sybil attacks on committee
    membership.
  \item \textbf{Stage 4}: Ablation study---remove AES, remove RS,
    remove Shamir individually to quantify each component's
    contribution.
  \item \textbf{TLA+ specification}: formalize the 6 invariants and
    model-check the protocol under all failure modes.
  \item \textbf{Production implementation}: Go node with BLS, KZG, and
    network transport.
  \item \textbf{Coq verification}: prove isomorphism between TLA+
    specification and production code.
\end{itemize}
